\section{Discussion, Boundaries, and Conclusion}
\label{sec:discussion}

\paragraph{What is complete.}
Pathwise history realization is complete for finite-horizon innovation-driven laws under the fixed coupling, M\"obius inversion is complete for lower-finite intervention responses, and the observation kernel is the complete indistinguishability gauge for a declared finite support.  Completeness concerns representation and identification under a fixed protocol; it does not provide a distribution-only minimal stochastic state or a universal predictive transition map across protocols.

\paragraph{What the interaction layer adds.}
P1 explains reduced-value interaction through inverse vertical stiffness, and P2 gives the path-space version.  The new readout theorem shows how that hidden response reaches an experimentally selected update, trace, or task observable.  Its extra readout and second-response terms explain why an actual optimizer response need not inherit the P1 negative-semidefinite sign.  The same observed pair effect can admit multiple hidden models, so the calculus does not identify geometry from responses alone.  Nonunique, nonsmooth, or hysteretic hidden responses require generalized derivatives or set-valued effects.

\paragraph{Projection and causality.}
A small residual proves proximity to a response class.  It does not allocate causal credit.  A passive top configuration leaves $d(|\Kcal|-1)$ supported coordinates ambiguous.  Controlled configurations, additional instrumentation, or structural restrictions are required for attribution.

\paragraph{The price of exact structure.}
Unrestricted Boolean attribution requires $2^m$ configurations, and its saturated Gaussian risk grows as $(3+2\sqrt2)^m/N$.  Scientific tractability comes from falsifiable structure: low interaction order, sparse hierarchy, graded interventions, or a target decision subspace.  Adaptive support selection requires selective-inference or multiplicity control beyond the fixed-support theorem.

\paragraph{Statistical boundary.}
Exact minimax, confidence, and chi-square conclusions assume a finite-dimensional linear experiment with known positive-definite covariance.  Estimated covariance, dependent traces, and adaptive designs need additional analysis.  Randomization calibration is exact only under a defensible invariant group action.

\paragraph{Empirical boundary.}
The digits experiment directly validates the P1 reduced-value curvature law and the P3 incidence, design, inference, and falsification layers under global strong convexity and additive interventions.  It does not validate the general readout-transfer law for an actual nonconvex optimizer trace.  The neural experiments establish response-class sensitivity under fixed protocols, but they lack factorial masks and hidden-response measurements.  Broad optimizer superiority would additionally require matched hyperparameter search, compute, throughput, memory, and modern workloads.

\paragraph{Research direction.}
The immediate theoretical target is to specialize the readout-transfer law to the full SPD and path-space realizations supplied by the companion papers.  The immediate empirical target is a nonconvex geometry--memory--operator campaign that measures update readouts, complete factorial masks, and the hidden first and second responses required by \cref{eq:readout-interaction-transfer}.

\paragraph{Conclusion.}
Causal optimizer interaction calculus separates what an optimizer can generate, what an intervention changes, and what an experiment identifies.  The universal layer covers innovation-driven finite-horizon rules under a fixed coupling.  P1 and P2 supply the hidden Schur structure; the new readout layer transports it to actual optimizer observables, and the incidence layer states exactly which resulting effects an experiment can recover.  This organization turns the four-paper geometry program into a concrete experimental theory without duplicating the companion structural theorems or forcing general optimizers into a gradient-flow model.
